%%%PAGE 272 rot=0 legibility=good summary=上部为竖直圆轨道+斜坡抛出的手写题（两问）及解答；中部贴有传送带甲乙的印刷题（4+4+4）及铅笔答案与手写解答；下半页空白
%%%BLOCK id=272-01 ch=04 sec=竖直面圆周·斜抛综合 kind=problem alt=06 cont=none y=0.03-0.30
%FIG p272_a 0.04 0.03 0.42 0.13 soft
\notefig[0.5]{figures/p272_a.png}
图中标注：$m=0.02\unit{kg}$，$V_0$（箭头），$R=0.4\unit{m}$，粗糙段$\mu=0.25$，斜坡端有$\alpha$。
\begin{problem}
(1) 滑块始终在轨道上运动，求$V_0$范围

(2) 滑块恰过D，求$E_{k\min}$时$X$与$\alpha$的关系，并求$X$最大值。
\begin{solution}
从\unclear{…}到A时速度为0
\[
mg(h-R)-\mu mgR=0-0\qquad h=0.5\unit{m}
\]
\[
mg(R-h)-\mu mgR=0-\frac{1}{2}mV_{01}^2 % 原稿“μ”写得像“4”
\]
\[
\therefore V_{01}=2\unit{m/s}\qquad 0\le V_{01}\le2\unit{m/s}
\]
\[
mg\,2R\cos\alpha=\frac{1}{2}mV_F^2
\]
\[
X=V_F\cos\alpha\cdot\frac{V_F\sin\alpha}{g}=\frac{V_F^2}{g}\sin\alpha\cos\alpha=1.6\cos^2\alpha\sin\alpha\le\frac{16\sqrt3}{45}\ \Rightarrow\ \tan\alpha=\frac{\sqrt2}{2}\ \text{时取等}
\]
\end{solution}
\end{problem}
%%%END
%%%BLOCK id=272-02 ch=06 sec=传送带·摩擦生热与电动机功率 kind=problem alt=03 cont=none y=0.30-0.79
% 贴纸印刷题；原稿“电动机的平均输出”处有铅笔划线
\begin{problem}[贴纸印刷题]
（4+4+4）如图所示，生产车间有两个相互垂直且等高的水平传送带甲和乙，甲的速度为$2v_0$。小工件离开甲前与甲的速度相同，并平稳地传到乙上，工件与乙之间的动摩擦因数为$\mu$。乙的宽度足够大，重力加速度为$g$。
\begin{itemize}
\item[(1)] 若乙的速度为$2v_0$，求工件在乙上侧向（垂直于乙的运动方向）滑过的距离$s$；
\item[(2)] 若乙的速度为$4v_0$，求工件在乙上刚停止侧向滑动时的速度大小$v$；
\item[(3)] 保持乙的速度$4v_0$不变，当工件在乙上刚停止滑动时，下一只工件恰好传到乙上，如此反复。若每个工件的质量均为$m$，除工件与传送带之间摩擦外，其他能量损耗均不计，求驱动乙的电动机的平均输出功率$\overline{P}$。
\end{itemize}

%FIG p272_b 0.02 0.48 0.27 0.625
\notefig[0.3]{figures/p272_b.png}

手写铅笔答案：(1)旁（箭头指向“距离$s$”）：
\[
S=\frac{2\sqrt2\,V_0^2}{\mu g}
\]
(2)旁（写在“速度大小$v$”下方）：$4V_0$。图中乙带右缘另有铅笔小标注\unclear{$\theta_1$}。

\begin{solution}
\begin{align*}
a_x&=\mu g\cos\theta\\
a_y&=\mu g\sin\theta\\
-2a_xS&=0-(2V_0)^2\\
2a_yy&=(4V_0)^2-0
\end{align*}
\[
t=\frac{4V_0}{a_y}\qquad y_1=4V_0t
\]
工件相对乙位移\ $L=\sqrt{x^2+(y_1-y)^2}$
\[
Q=\mu mgL\qquad W=\frac{1}{2}m(4V_0)^2-\frac{1}{2}mV_0^2+Q
\] % CHECK: 初动能项原稿为½mV₀²，工件进入乙时速度为2v₀，应为½m(2V₀)²？
\[
P=\frac{W}{t}\qquad \overline{P}=\frac{8\sqrt5\,\mu mgV_0}{5}
\]
\end{solution}
\end{problem}
%%%END
%%%PAGEEND 272
